\documentclass[11pt,a4paper]{article}

\usepackage[utf8]{inputenc}
\usepackage[T1]{fontenc}
\usepackage[french]{babel}
\usepackage{lmodern}
\usepackage[margin=2.5cm]{geometry}
\usepackage{amsmath,amssymb,amsthm}
\usepackage{booktabs}
\usepackage{pgfplots}
\pgfplotsset{compat=1.18}
\usepackage[hidelinks]{hyperref}

\newtheorem{proposition}{Proposition}
\newtheorem{lemme}{Lemme}
\newtheorem{corollaire}{Corollaire}
\theoremstyle{definition}
\newtheorem{definition}{Définition}
\theoremstyle{remark}
\newtheorem*{remarque}{Remarque}

\newcommand{\HIT}{\textsc{hit}}
\newcommand{\STAND}{\textsc{stand}}
\newcommand{\BUST}{\textsc{bust}}
\newcommand{\QH}{Q_{\HIT}}
\newcommand{\VS}{V_{\STAND}}
\newcommand{\E}{\mathbb{E}}
\newcommand{\1}{\mathbf{1}}
\newcommand{\KL}{\mathrm{KL}}
\newcommand{\TV}{\mathrm{TV}}
\newcommand{\R}{\mathcal{R}}

\title{Blackjack AI --- Fonctionnement mathématique du modèle 1\\
\large Prédiction probabiliste du croupier et effet de son erreur sur l'agent hybride}
\author{}
\date{}

\begin{document}
\maketitle

\tableofcontents

\bigskip

Ce document complète \emph{Fonctionnement mathématique du modèle 2} (fichier
\texttt{docs/modele2.pdf}), qui décrit le processus de décision markovien, le Q-learning
hybride et la propagation d'une erreur du modèle 1 vers la politique. Nous en reprenons
les notations : $P(c)$ désigne la loi d'une carte, $d$ la carte visible du croupier,
$D\in\mathcal{D}=\{17,18,19,20,21,\BUST\}$ son résultat final et $R(p,D)$ le gain du
joueur qui reste avec le total $p$. Le présent document traite de la partie statistique :
ce que le modèle 1 estime, comment il l'estime, avec quelle précision en fonction du
nombre de mains simulées, et comment cette précision se traduit en qualité de décision.
Toutes les valeurs numériques citées proviennent des fichiers de résultats du dépôt ou
ont été recalculées à partir des modèles livrés.

%======================================================================
\section{Le problème d'apprentissage}
%======================================================================

\subsection{Données}

On note $\mathcal{C}=\{1,\dots,10\}$ l'ensemble des cartes visibles possibles et
$K=|\mathcal{D}|=6$ le nombre de résultats du croupier. Une simulation tire la carte
visible $d$ selon $P$, puis la carte cachée et les cartes suivantes, et applique la règle
du croupier jusqu'à son arrêt. On obtient un exemple $(d_i,D_i)$. Les $n$ exemples d'un jeu
de données sont indépendants et de même loi :
\[
d_i\sim P,\qquad D_i\mid d_i\sim\pi_{d_i},\qquad\text{où }\pi_d(k)=P(D=k\mid d).
\]
On note $n_d$ le nombre d'exemples de carte visible $d$, $n_{dk}$ le nombre de ceux qui
se terminent par le résultat $k$, et $f_d(k)=n_{dk}/n_d$ la fréquence observée
correspondante.

\subsection{La carte visible suffit}

\begin{proposition}\label{prop:suffisance}
Dans un paquet infini, pour tout état $s=(p,d,a)$ dans lequel le joueur reste,
\[
P(D=k\mid s)=\pi_d(k)\qquad\text{pour tout }k\in\mathcal{D}.
\]
\end{proposition}

\begin{proof}
Toutes les cartes sont des tirages indépendants de loi $P$. Le résultat $D$ est une
fonction de la carte visible, de la carte cachée et des cartes que tire ensuite le
croupier, car sa règle d'arrêt ne dépend que de sa propre main. Les cartes du joueur sont
d'autres tirages, indépendants des précédents. Conditionnellement à $d$, le résultat $D$
est donc indépendant de $(p,a)$.
\end{proof}

Cette propriété justifie de simuler le croupier séparément du joueur et de ne donner au
modèle que la carte visible. Elle cesserait d'être vraie avec un paquet fini, où les
cartes déjà distribuées au joueur modifient la composition du paquet.

\subsection{Ce qu'il faut estimer}

Le modèle 1 doit estimer dix lois de probabilité $\pi_1,\dots,\pi_{10}$, chacune élément
du simplexe
\[
\Delta_K=\Big\{q\in[0,1]^K : \sum_{k}q(k)=1\Big\}.
\]
Chaque loi a $K-1=5$ degrés de liberté, soit $50$ paramètres libres au total. Il s'agit
d'un problème d'estimation de faible dimension, pour lequel on dispose de dizaines de
milliers d'exemples : on s'attend donc à une erreur faible, que la
section~\ref{sec:taille} quantifie.

%======================================================================
\section{Évaluer une prédiction probabiliste}
%======================================================================

Une prédiction est une famille $q=(q_1,\dots,q_{10})$ de lois de $\Delta_K$. On l'évalue
par des \emph{règles de score} qui comparent la loi annoncée au résultat observé.

\subsection{Log-loss, entropie et divergence de Kullback-Leibler}

\begin{definition}
La log-loss d'une prédiction $q$ sur l'exemple $(d,D)$ vaut $-\log q_d(D)$, en nats.
Son espérance sous la vraie loi des données est le risque
\[
\R(q)=\E\big[-\log q_d(D)\big]=-\sum_{d\in\mathcal{C}}P(d)\sum_{k\in\mathcal{D}}\pi_d(k)\log q_d(k).
\]
\end{definition}

\begin{proposition}[Décomposition du risque]\label{prop:decomposition}
Pour toute prédiction $q$ dont les probabilités sont strictement positives,
\[
\R(q)=H(D\mid d)+\sum_{d\in\mathcal{C}}P(d)\,\KL(\pi_d\,\|\,q_d),
\]
où $H(D\mid d)=-\sum_d P(d)\sum_k\pi_d(k)\log\pi_d(k)$ est l'entropie conditionnelle du
résultat et $\KL(\pi\,\|\,q)=\sum_k\pi(k)\log\big(\pi(k)/q(k)\big)$ la divergence de
Kullback-Leibler.
\end{proposition}

\begin{proof}
On écrit $-\log q_d(k)=-\log\pi_d(k)+\log\big(\pi_d(k)/q_d(k)\big)$, puis on prend
l'espérance sous $\pi_d$ et sous $P$.
\end{proof}

La divergence de Kullback-Leibler est positive et ne s'annule que si les deux lois sont
égales (inégalité de Gibbs). On en déduit trois conséquences.
\begin{enumerate}
  \item Le risque est minimal pour $q=\pi$ et seulement pour elle : la log-loss
  récompense les probabilités exactes, ce que l'on exprime en disant que c'est une règle
  de score \emph{strictement propre}~\cite{gneiting}.
  \item La valeur minimale est l'entropie conditionnelle, qui vaut ici $1{,}6481$ nat
  par main. Aucun modèle ne peut faire mieux en moyenne : cette part de la perte
  traduit le hasard des cartes, pas une imperfection du modèle.
  \item Le second terme, appelé \emph{excès de log-loss}, mesure seulement l'erreur du
  modèle. Comme $\pi$ est connue exactement (section 2.1 du document du modèle 2), cet
  excès se calcule sans aucun jeu de test. Pour le modèle livré, il vaut $0{,}000429$.
\end{enumerate}

\begin{remarque}
La log-loss mesurée sur un jeu de test fini estime $\R(q)$ avec un bruit important.
L'écart-type d'une perte individuelle vaut environ $0{,}55$, donc celui d'une moyenne
sur $12\,000$ mains vaut environ $0{,}005$, soit douze fois l'excès à mesurer. C'est
pourquoi la loi exacte elle-même peut obtenir une log-loss de test légèrement moins bonne
qu'un modèle estimé ($1{,}651709$ contre $1{,}651587$), et pourquoi les comparaisons
utilisent l'excès exact ou des différences appariées sur les mêmes mains.
\end{remarque}

\subsection{Score de Brier}

\begin{proposition}\label{prop:brier}
Le score de Brier multiclasse $B(q)=\E\big[\sum_k\big(q_d(k)-\1[D=k]\big)^2\big]$ se
décompose en
\[
B(q)=\sum_{d}P(d)\Big(\|q_d-\pi_d\|^2+1-\|\pi_d\|^2\Big).
\]
\end{proposition}

\begin{proof}
Pour $d$ fixé et $Y_k=\1[D=k]$, on a $\E[Y_k]=\pi_d(k)$ et
$\E\big[(q_d(k)-Y_k)^2\big]=(q_d(k)-\pi_d(k))^2+\pi_d(k)(1-\pi_d(k))$. On somme sur $k$
en utilisant $\sum_k\pi_d(k)=1$.
\end{proof}

Le score de Brier est donc lui aussi strictement propre. Sa part irréductible,
$\sum_d P(d)\big(1-\|\pi_d\|^2\big)$, vaut $0{,}7768$ : comme pour la log-loss, la
quasi-totalité du score observé ($0{,}7785$ sur le test) provient du hasard des cartes.

\subsection{Exactitude}

Un prédicteur qui annonce toujours le résultat le plus probable a pour exactitude
\[
\sum_{d}P(d)\,\pi_d\big(\arg\max_k q_d(k)\big)\le\sum_d P(d)\max_k\pi_d(k)=36{,}70\,\%.
\]
La borne est atteinte dès que $q_d$ et $\pi_d$ ont le même mode pour chaque carte, quelle
que soit la précision des autres probabilités. C'est le cas pour la régression, pour les
fréquences et pour la loi exacte, qui obtiennent toutes trois $36{,}49\,\%$ sur le même
jeu de test. L'exactitude ne permet donc pas de distinguer ces estimateurs ; elle
n'évalue pas non plus les probabilités dont l'agent hybride a besoin.

%======================================================================
\section{La régression logistique multinomiale}
%======================================================================

\subsection{Le modèle}

La carte visible est représentée par le vecteur indicateur $x=e_d\in\{0,1\}^{10}$
(encodage \emph{one-hot}). Le modèle associe à chaque résultat $k$ un vecteur de poids
$w_k\in\mathbb{R}^{10}$ et un biais $b_k\in\mathbb{R}$. Comme $w_k^\top e_d=w_{k,d}$, le
score du résultat $k$ pour la carte $d$ vaut $z_k(d)=w_{k,d}+b_k$, et la probabilité
prédite est donnée par la fonction \emph{softmax}~\cite{hastie} :
\[
q_d(k)=\frac{\exp\big(w_{k,d}+b_k\big)}{\sum_{j\in\mathcal{D}}\exp\big(w_{j,d}+b_j\big)}.
\]
Les paramètres forment une matrice $W\in\mathbb{R}^{6\times 10}$ et un vecteur
$b\in\mathbb{R}^6$, soit 66 nombres pour 50 degrés de liberté : ajouter une même
constante à tous les scores d'une carte ne change pas ses probabilités. La
régularisation, décrite ci-dessous, sélectionne une représentation parmi toutes
celles qui donnent les mêmes probabilités.

\subsection{Estimation}

Avec le solveur L-BFGS, scikit-learn minimise
\[
\mathcal{L}_C(W,b)=\frac12\|W\|_F^2+C\sum_{i=1}^n-\log q_{d_i}(D_i),
\]
où $\|W\|_F^2=\sum_{k,d}w_{k,d}^2$ et où les biais ne sont pas pénalisés. Diviser par
$Cn$ montre qu'il s'agit de minimiser la log-loss moyenne d'entraînement plus la pénalité
$\|W\|_F^2/(2Cn)$ : plus $C$ est petit, plus la pénalité est forte. De manière
équivalente, $(W,b)$ est l'estimateur du maximum a posteriori lorsque chaque poids suit a
priori une loi normale centrée de variance $C$. La forme exacte de cet objectif a été
vérifiée numériquement à travers sa condition d'optimalité
(proposition~\ref{prop:regularisation}).

\begin{proposition}[Convexité et gradient]\label{prop:gradient}
La fonction $\mathcal{L}_C$ est convexe, et strictement convexe en $W$. Ses dérivées
partielles valent
\[
\frac{\partial\mathcal{L}_C}{\partial w_{k,d}}=w_{k,d}+C\,n_d\big(q_d(k)-f_d(k)\big),
\qquad
\frac{\partial\mathcal{L}_C}{\partial b_k}=C\sum_{d}n_d\big(q_d(k)-f_d(k)\big).
\]
\end{proposition}

\begin{proof}
La perte d'un exemple s'écrit $-z_{D_i}(d_i)+\log\sum_j\exp z_j(d_i)$. La fonction
$\log\sum\exp$ est convexe et les scores sont affines en $(W,b)$, donc la perte est
convexe ; le terme $\frac12\|W\|_F^2$ est strictement convexe en $W$. La dérivée de la
perte d'un exemple par rapport à $z_k$ vaut $q_{d_i}(k)-\1[D_i=k]$. Seuls les exemples de
carte $d$ font intervenir $w_{k,d}$ ; leur somme donne
$n_d\,q_d(k)-n_{dk}=n_d\big(q_d(k)-f_d(k)\big)$. Le biais $b_k$ intervient dans tous les
exemples, d'où la somme sur $d$.
\end{proof}

Annuler ce gradient conduit à des équations non linéaires, car $q_d$ dépend de $W$ et
de $b$ à travers la softmax. Il n'existe pas de formule explicite pour $C$ fini, et
l'optimum est calculé numériquement par L-BFGS, une méthode de quasi-Newton qui
approche la courbure de l'objectif à partir des gradients
successifs~\cite{nocedal}. La convexité garantit que le minimum trouvé est global.

\begin{remarque}[Ce n'est pas une méthode des moindres carrés]
Les moindres carrés ajustent une fonction linéaire en minimisant une somme d'écarts
quadratiques ; leurs prédictions ne sont pas contraintes à former des probabilités. La
régression logistique maximise au contraire la vraisemblance d'un modèle de
probabilités, ce qui revient à minimiser la log-loss. Les deux approches coïncident
toutefois dans le cas particulier étudié ici : avec une seule variable catégorielle, une
régression linéaire des indicatrices $\1[D=k]$ sur l'encodage one-hot donne elle aussi
les fréquences $f_d(k)$, qui sont les moyennes de ces indicatrices dans chaque catégorie.
Les deux méthodes ne divergent qu'en présence de plusieurs variables explicatives.
\end{remarque}

%======================================================================
\section{Le cas d'une seule variable catégorielle}
%======================================================================

\subsection{Sans régularisation, la régression redonne les fréquences}

\begin{proposition}[Modèle saturé]\label{prop:sature}
Supposons $n_{dk}>0$ pour toute carte $d$ et tout résultat $k$. Alors la log-loss
d'entraînement $\sum_i-\log q_{d_i}(D_i)$ atteint son minimum exactement lorsque
$q_d=f_d$ pour toute carte $d$.
\end{proposition}

\begin{proof}
En regroupant les exemples par carte, puis en appliquant la
proposition~\ref{prop:decomposition} à la loi empirique,
\[
\sum_{i=1}^n-\log q_{d_i}(D_i)=\sum_{d}n_d\Big(H(f_d)+\KL(f_d\,\|\,q_d)\Big).
\]
Le premier terme ne dépend pas du modèle et le second est minimal, égal à zéro, si et
seulement si $q_d=f_d$ pour chaque $d$. Cette valeur est atteignable : il suffit de
prendre $b=0$ et $w_{k,d}=\log f_d(k)$, ce qui est possible car toutes les fréquences
sont strictement positives.
\end{proof}

Le modèle est dit \emph{saturé} : il peut représenter n'importe quelle famille de dix
lois strictement positives, et rien ne relie les estimations des différentes cartes. Si
un résultat n'était jamais observé pour une carte, la fréquence correspondante serait
nulle, valeur qu'une softmax ne peut pas atteindre, et les poids divergeraient. Dans le
jeu d'entraînement livré, le plus petit effectif $n_{dk}$ vaut $278$ : la situation ne
se produit pas.

\subsection{L'effet exact de la régularisation}

\begin{proposition}\label{prop:regularisation}
Au minimum de $\mathcal{L}_C$, pour toute carte $d$ et tout résultat $k$,
\[
q_d(k)-f_d(k)=-\frac{w_{k,d}}{C\,n_d}
\qquad\text{et}\qquad
\sum_{d\in\mathcal{C}}w_{k,d}=0.
\]
En particulier, $\sum_d n_d\,q_d(k)=\sum_d n_{dk}$ : le modèle reproduit exactement la
fréquence globale de chaque résultat dans l'échantillon d'entraînement.
\end{proposition}

\begin{proof}
On annule les dérivées de la proposition~\ref{prop:gradient}. La dérivée par rapport à
$w_{k,d}$ donne la première égalité. La dérivée par rapport à $b_k$ donne
$\sum_d n_d\big(q_d(k)-f_d(k)\big)=0$ ; en y remplaçant chaque terme par
$-w_{k,d}/C$, on obtient la seconde.
\end{proof}

Cette proposition décrit entièrement l'effet de la régularisation. Lorsque les poids
sont nuls, toutes les cartes reçoivent la même loi, $\mathrm{softmax}(b)$ ; la
pénalité attire donc chaque $q_d$ vers une loi commune à toutes les cartes. L'écart aux
fréquences est inversement proportionnel à $C\,n_d$ : il est plus marqué pour les cartes
rares et pour les petites valeurs de $C$.

Pour le modèle livré, l'identité a été vérifiée sur les soixante probabilités avec un
écart maximal de $1{,}5\times10^{-7}$, dû à la tolérance de l'optimiseur. Avec
$C=0{,}1$ et $n_d$ compris entre $4\,546$ et $18\,374$, on a $C\,n_d\ge 454$, et le plus
grand écart entre la régression et les fréquences vaut $0{,}0026$. Cette proposition
explique trois observations du projet.
\begin{enumerate}
  \item \emph{Le choix de $C$ importe peu.} Pour $C\ge 0{,}1$, les écarts aux
  fréquences sont trop petits pour modifier la log-loss de validation au-delà de la
  sixième décimale. Pour $C=0{,}01$, ils sont dix fois plus grands et la log-loss de
  validation se dégrade nettement ($1{,}641715$ contre $1{,}641194$).
  \item \emph{La régression et les fréquences ont la même qualité.} Leur excès de
  log-loss diffère de $0{,}00002$ à $60\,000$ mains, alors que l'excès lui-même vaut
  $0{,}0004$.
  \item \emph{Les deux agents hybrides prennent presque toujours les mêmes décisions.}
  D'après le lemme sur la variation totale du document du modèle 2, les valeurs de
  \STAND{} calculées avec la régression et avec les fréquences diffèrent d'au plus
  $2\max_d\TV(q_d,f_d)=0{,}0053$. Les deux politiques ne peuvent donc différer que dans
  des états où les valeurs de \HIT{} et de \STAND{} sont presque égales, ce qui est
  cohérent avec les gains identiques observés sur 18 des 20 graines de l'ablation.
\end{enumerate}

\begin{remarque}[Quand la régression serait-elle utile ?]
L'intérêt d'un modèle paramétrique apparaît lorsqu'il n'est pas saturé, c'est-à-dire
lorsque le nombre de situations possibles est grand devant le nombre de paramètres. Avec
un paquet fini, par exemple, la loi du croupier dépendrait aussi des cartes déjà
sorties ; les combinaisons possibles seraient trop nombreuses pour être estimées
séparément, et un modèle qui partage l'information entre situations voisines deviendrait
nécessaire.
\end{remarque}

%======================================================================
\section{Erreur d'estimation et quantité de données}\label{sec:taille}
%======================================================================

\subsection{Comportement asymptotique}

\begin{proposition}\label{prop:asymptotique}
Soit $f$ le vecteur des fréquences de $m$ tirages indépendants de loi
$\pi\in\Delta_K$, avec $\pi(k)>0$ pour tout $k$. Lorsque toutes les fréquences sont
strictement positives,
\[
\KL(\pi\,\|\,f)=\frac12\sum_{k}\frac{\big(f(k)-\pi(k)\big)^2}{\pi(k)}+O_P\big(m^{-3/2}\big),
\qquad
\E\bigg[\frac12\sum_k\frac{\big(f(k)-\pi(k)\big)^2}{\pi(k)}\bigg]=\frac{K-1}{2m}.
\]
\end{proposition}

\begin{proof}
On pose $\delta_k=f(k)-\pi(k)$ ; on a $\sum_k\delta_k=0$ et $\delta_k=O_P(m^{-1/2})$.
Le développement $-\log(1+u)=-u+u^2/2+O(u^3)$ appliqué à $u=\delta_k/\pi(k)$ donne
\[
\KL(\pi\,\|\,f)=-\sum_k\pi(k)\log\Big(1+\frac{\delta_k}{\pi(k)}\Big)
=-\sum_k\delta_k+\frac12\sum_k\frac{\delta_k^2}{\pi(k)}+O_P\big(m^{-3/2}\big),
\]
et le terme linéaire est nul. Pour la loi multinomiale,
$\E[\delta_k^2]=\pi(k)\big(1-\pi(k)\big)/m$, d'où
$\E\big[\frac12\sum_k\delta_k^2/\pi(k)\big]=\sum_k\big(1-\pi(k)\big)/(2m)=(K-1)/(2m)$.
\end{proof}

La probabilité qu'une fréquence soit nulle décroît exponentiellement avec $m$. Pour
les tailles étudiées ici, elle est négligeable : avec $5\,000$ mains, une carte est
observée environ $385$ fois et le résultat le moins probable a une probabilité d'environ
$0{,}06$, si bien qu'il manque avec une probabilité de l'ordre de $10^{-10}$.

\begin{corollaire}\label{cor:25n}
Pour $n$ mains d'entraînement, l'excès moyen de log-loss de l'estimateur par fréquences
vaut, à l'ordre principal,
\[
\sum_{d\in\mathcal{C}}P(d)\,\frac{K-1}{2\,n\,P(d)}=\frac{|\mathcal{C}|\,(K-1)}{2n}=\frac{25}{n}.
\]
\end{corollaire}

\begin{proof}
La carte $d$ est observée environ $n_d\approx n\,P(d)$ fois ; on applique la
proposition~\ref{prop:asymptotique} à chaque carte, puis on pondère par $P(d)$.
\end{proof}

Ce résultat ne dépend pas de la loi des cartes visibles : une carte rare est estimée
moins précisément, mais elle pèse moins dans le risque, et les deux effets se
compensent exactement. L'erreur diminue comme $1/n$ : il faut multiplier le nombre de
mains par dix pour diviser l'excès par dix.

Pour la régression régularisée, la proposition~\ref{prop:regularisation} montre que
$q_d-f_d=O(1/n_d)$. Cette correction ne modifie le développement de la
proposition~\ref{prop:asymptotique} qu'aux ordres supérieurs, et l'excès de la
régression suit donc la même loi $25/n$ à l'ordre principal. Pour de petits
échantillons, la correction n'est plus négligeable : avec $5\,000$ mains, $C\,n_d$
vaut environ $38$ au lieu de $460$, et l'attraction vers la loi commune réduit un peu
la variance de l'estimation. C'est cohérent avec le léger avantage, non significatif, de
la régression observé à cette taille.

\subsection{Comparaison avec les mesures}

L'étude de la quantité de données répète l'entraînement sur cinq jeux indépendants pour
quatre tailles. Le tableau~\ref{tab:25n} et la figure~\ref{fig:25n} comparent les excès
moyens mesurés à la prédiction du corollaire~\ref{cor:25n}.

\begin{table}[h]
\centering
\begin{tabular}{rrrr}
\toprule
Mains d'entraînement $n$ & Prédiction $25/n$ & Fréquences (mesure) & Régression (mesure) \\
\midrule
$5\,000$   & $0{,}005000$ & $0{,}005007$ & $0{,}004507$ \\
$20\,000$  & $0{,}001250$ & $0{,}001310$ & $0{,}001257$ \\
$60\,000$  & $0{,}000417$ & $0{,}000434$ & $0{,}000424$ \\
$200\,000$ & $0{,}000125$ & $0{,}000119$ & $0{,}000129$ \\
\bottomrule
\end{tabular}
\caption{Excès de log-loss moyen sur cinq jeux indépendants, comparé à la prédiction
asymptotique.}\label{tab:25n}
\end{table}

\begin{figure}[h]
\centering
\begin{tikzpicture}
\begin{loglogaxis}[
  width=0.78\textwidth, height=6.2cm,
  xlabel={Mains d'entraînement $n$},
  ylabel={Excès de log-loss (nat par main)},
  legend pos=north east, legend cell align=left,
  grid=major, grid style={gray!25},
]
\addplot[domain=4000:250000, samples=2, thick, black] {25/x};
\addlegendentry{Prédiction $25/n$}
\addplot[only marks, mark=*, blue] coordinates {
  (5000,0.004507) (20000,0.001257) (60000,0.000424) (200000,0.000129)};
\addlegendentry{Régression logistique}
\addplot[only marks, mark=square*, orange] coordinates {
  (5000,0.005007) (20000,0.001310) (60000,0.000434) (200000,0.000119)};
\addlegendentry{Fréquences}
\end{loglogaxis}
\end{tikzpicture}
\caption{Excès de log-loss moyen en fonction du nombre de mains, en échelle
logarithmique. La droite de pente $-1$ est la prédiction du
corollaire~\ref{cor:25n}.}\label{fig:25n}
\end{figure}

Sur deux ordres de grandeur de taille d'échantillon, les fréquences suivent la prédiction
à moins de 5\,\% près. La régression la suit de la même manière, sauf à $5\,000$ mains où
elle fait 10\,\% mieux, pour la raison donnée ci-dessus. L'erreur résiduelle du modèle 1
s'explique donc entièrement par la taille finie des données : le modèle saturé n'a pas
de biais d'approximation, et il n'y a pas de surapprentissage à corriger.

\begin{remarque}[Dispersion d'une carte à l'autre]
La quantité $2m\cdot\frac12\sum_k\delta_k^2/\pi(k)$ est la statistique du $\chi^2$ de
Pearson. Asymptotiquement, $2\,n_d\KL(\pi_d\,\|\,f_d)$ suit donc approximativement une loi
du $\chi^2$ à $K-1=5$ degrés de liberté, de moyenne 5 et d'écart-type $\sqrt{10}\approx3{,}2$.
Dans le jeu d'entraînement livré, cette quantité vaut en moyenne $5{,}4$ sur les dix
cartes. Elle atteint $13{,}7$ pour la carte 6, une valeur qu'une loi du $\chi^2$ à cinq
degrés de liberté dépasse avec une probabilité de $1{,}8\,\%$. Sur dix cartes, la
probabilité qu'au moins une atteigne une telle valeur est de $16\,\%$ : c'est une
fluctuation ordinaire. Elle explique toutefois pourquoi la carte 6 concentre les plus
grandes erreurs du modèle livré (section~\ref{sec:decision}).
\end{remarque}

%======================================================================
\section{De l'erreur des probabilités à l'erreur de décision}\label{sec:decision}
%======================================================================

\subsection{Erreur sur la valeur de rester}

\begin{proposition}\label{prop:pinsker}
Pour tout total $p$ et toute carte $d$,
\[
\big|\hat V_{\STAND}(p,d)-\VS(p,d)\big|\le2\,\TV(q_d,\pi_d)\le\sqrt{2\,\KL(\pi_d\,\|\,q_d)}.
\]
De plus, pour $p\le16$, $\hat V_{\STAND}(p,d)-\VS(p,d)=2\big(q_d(\BUST)-\pi_d(\BUST)\big)$.
\end{proposition}

\begin{proof}
La première inégalité est le lemme sur la variation totale du document du modèle 2
(section 6). La seconde est l'inégalité de Pinsker,
$\TV(\pi,q)\le\sqrt{\KL(\pi\,\|\,q)/2}$~\cite{cover}. Enfin, pour $p\le16$, le joueur ne
gagne que si le croupier dépasse 21, donc $\VS(p,d)=2\pi_d(\BUST)-1$, et de même pour
$\hat V_{\STAND}$ (section 2 du document du modèle 2).
\end{proof}

En combinant cette proposition avec la proposition~\ref{prop:asymptotique}, on obtient
l'ordre de grandeur typique de l'erreur sur la valeur de rester :
\[
\big|\hat V_{\STAND}-\VS\big|\lesssim\sqrt{2\cdot\frac{K-1}{2\,n_d}}=\sqrt{\frac{5}{n\,P(d)}}.
\]
Avec $n=60\,000$, cela donne environ $0{,}033$ pour une carte de 1 à 9 et $0{,}017$ pour
la carte 10. L'erreur ne décroît qu'en $1/\sqrt{n}$ : la diviser par deux demande quatre
fois plus de mains.

Pour le modèle livré, la plus grande distance en variation totale vaut
$\varepsilon_{TV}=0{,}0229$ et est atteinte pour la carte 6, d'où la borne
$\eta=2\varepsilon_{TV}=0{,}0457$. L'erreur maximale effectivement observée sur
$\VS$ vaut $0{,}0387$, dans l'état $(17,6,0)$. Pour les totaux inférieurs ou égaux à 16,
l'erreur provient uniquement de la probabilité de dépassement : pour la carte 6, le
modèle annonce $0{,}4041$ au lieu de $0{,}4232$, ce qui donne une erreur de
$2\times0{,}0191=0{,}0382$.

\subsection{Conséquences pour la politique}

Le document du modèle 2 (section 6) établit trois résultats qui relient $\eta$ à la
politique hybride limite, c'est-à-dire à celle qu'obtiendrait le Q-learning hybride
après un nombre infini d'épisodes :
\begin{enumerate}
  \item la valeur de \HIT{} limite vérifie $|\hat Q_{\HIT}-\QH^*|\le\eta$ dans tout état ;
  \item la décision ne peut différer de la décision optimale que dans les états
  \emph{ambigus}, où $|\QH^*-\VS|\le2\eta$ ;
  \item la perte de gain est au plus $36\,\eta$.
\end{enumerate}
Le tableau~\ref{tab:bornes} confronte ces bornes aux valeurs calculées pour le modèle
livré, en résolvant les équations de Bellman avec $\hat V_{\STAND}$.

\begin{table}[h]
\centering
\begin{tabular}{lrr}
\toprule
Quantité & Borne théorique & Valeur calculée \\
\midrule
Erreur maximale sur $\VS$ & $\eta=0{,}0457$ & $0{,}0387$ \\
Erreur maximale sur $\QH$ & $\eta=0{,}0457$ & $0{,}0370$ \\
États où la décision peut différer & 27 états ambigus sur 280 & 1 état sur 280 \\
Perte de gain espéré par partie & $36\,\eta=1{,}65$ & $0{,}000022$ \\
\bottomrule
\end{tabular}
\caption{Bornes du document du modèle 2 et valeurs obtenues avec le modèle 1 livré.}\label{tab:bornes}
\end{table}

Les bornes sur les valeurs sont proches des erreurs observées, mais les bornes sur la
politique sont très pessimistes. Le seul état où la décision change est $(12,4,0)$ : la
politique optimale y reste, car $\VS=-0{,}2111$ dépasse $\QH^*=-0{,}2135$ de seulement
$0{,}0025$, alors que la politique hybride limite tire, car
$\hat Q_{\HIT}=-0{,}2239$ dépasse $\hat V_{\STAND}=-0{,}2273$. La perte qui en résulte,
$0{,}000022$ par partie, est négligeable devant l'écart de $0{,}00065$ qui sépare l'agent
livré de l'optimum après $100\,000$ épisodes.

\begin{remarque}[Pourquoi les erreurs se compensent en partie]
La décision dépend de la différence $\hat Q_{\HIT}-\hat V_{\STAND}$, et non de chaque
valeur séparément. Or $\hat Q_{\HIT}(s)$ est calculée à partir des valeurs
$\hat V_{\STAND}$ des états suivants, qui partagent la même carte visible et donc la même
erreur de probabilité. Les deux erreurs vont dans le même sens et s'annulent en partie.
Dans l'état $(12,4,0)$, la valeur de rester est sous-estimée de $0{,}0163$ et celle de
tirer de $0{,}0104$ : leur différence ne varie que de $0{,}0059$, alors que la borne
autorise une variation de $2\eta=0{,}0915$.
\end{remarque}

%======================================================================
\section{Protocole statistique}
%======================================================================

\paragraph{Choix de la régularisation.} Pour chaque valeur candidate de $C$, la log-loss
moyenne sur le jeu de validation est un estimateur sans biais du risque $\R(q_C)$ du
modèle entraîné. Choisir la valeur qui minimise ces estimations rend en revanche le score
retenu optimiste, car on sélectionne en partie le bruit de la validation. C'est pourquoi
la qualité finale est mesurée sur un jeu de test distinct, généré après ce choix.

\paragraph{Comparaison appariée sur le test.} Pour comparer deux estimateurs $q$ et
$q'$, on calcule sur chaque main de test la différence
$\Delta_i=-\log q_{d_i}(D_i)+\log q'_{d_i}(D_i)$. Le bruit dû aux cartes, commun aux
deux estimateurs, disparaît en grande partie de cette différence. L'intervalle
$\bar\Delta\pm1{,}96\,s_\Delta/\sqrt{m}$, où $s_\Delta$ est l'écart-type empirique des
$m=12\,000$ différences, est un intervalle de confiance approximatif à $95\,\%$ fondé sur
le théorème central limite. Entre la régression et les fréquences, il vaut
$[-0{,}000125\,;\,+0{,}000014]$ et contient zéro.

\paragraph{Répétitions indépendantes.} Pour l'étude de la quantité de données, l'unité
statistique est le jeu de données complet. Avec cinq jeux, l'intervalle de confiance de
la moyenne utilise le quantile de Student à quatre degrés de liberté, $t_{0{,}975}=2{,}776$,
au lieu de $1{,}96$. Les tailles d'entraînement d'une même répétition sont des préfixes du
même jeu : les statistiques sont donc calculées entre répétitions, à taille fixée.

\paragraph{Calibration par issue.} La figure de calibration du dépôt compare chaque
probabilité prédite à la fréquence observée sur le test, avec un intervalle de
Wilson~\cite{wilson}. Pour une proportion observée $\hat p$ sur $m$ mains et
$z=1{,}96$, cet intervalle vaut
\[
\frac{\hat p+\frac{z^2}{2m}\pm z\sqrt{\frac{\hat p(1-\hat p)}{m}+\frac{z^2}{4m^2}}}{1+\frac{z^2}{m}}.
\]
Il reste fiable pour les petites proportions, contrairement à l'intervalle
$\hat p\pm z\sqrt{\hat p(1-\hat p)/m}$. Les soixante intervalles sont calculés point par
point : ils ne forment pas une bande de confiance simultanée, et quelques points hors de
leur intervalle sont attendus même pour un modèle parfait.

%======================================================================
\section{Synthèse}
%======================================================================

Le modèle 1 estime dix lois de probabilité à partir d'une seule variable catégorielle.
Dans ce cadre, la régression logistique multinomiale est un modèle saturé : sans
régularisation, elle redonne exactement les fréquences observées, et la régularisation
ne l'en écarte que d'une quantité inversement proportionnelle à $C\,n_d$. Son erreur suit
la loi asymptotique $25/n$ de l'estimation par fréquences, ce que confirment les mesures
sur deux ordres de grandeur. Cette erreur se transmet à la valeur de rester avec une
amplitude de l'ordre de $\sqrt{5/(n\,P(d))}$, soit quelques centièmes pour
$60\,000$ mains. Elle ne modifie pourtant la politique limite que dans un seul des
280 états, pour une perte de gain négligeable. Ces résultats expliquent pourquoi
l'ablation ne distingue pas la régression des fréquences, et pourquoi l'avantage de
l'agent hybride tient à l'information apportée par les mains simulées plutôt qu'à la
forme du classifieur.

\begin{thebibliography}{9}
\bibitem{cover} T.~M. Cover et J.~A. Thomas, \emph{Elements of Information Theory},
2\textsuperscript{e} édition, Wiley, 2006.
\bibitem{gneiting} T.~Gneiting et A.~E. Raftery, « Strictly proper scoring rules,
prediction, and estimation », \emph{Journal of the American Statistical Association},
102(477), 2007, p.~359--378.
\bibitem{hastie} T.~Hastie, R.~Tibshirani et J.~Friedman, \emph{The Elements of
Statistical Learning}, 2\textsuperscript{e} édition, Springer, 2009, chapitre~4.
\bibitem{nocedal} J.~Nocedal et S.~J. Wright, \emph{Numerical Optimization},
2\textsuperscript{e} édition, Springer, 2006, chapitre~7.
\bibitem{wilson} E.~B. Wilson, « Probable inference, the law of succession, and
statistical inference », \emph{Journal of the American Statistical Association},
22(158), 1927, p.~209--212.
\end{thebibliography}

\end{document}
